\chapter{Conclusion}

A substantial foundation has been established through an extensive literature review, a clear problem definition, a detailed methodology, and a thorough understanding of the dataset. Emphasis has been placed on recognizing the complexity and breadth of ICD coding tasks, with particular attention to the underrepresentation of rare conditions. Although progress has been steady, the generation of synthetic data has not yet been executed, and it remains a pivotal step for enhancing the coverage and accuracy of ICD assignments.

Moving forward, the methodology’s execution continues to be refined, with plans to employ knowledge-driven data augmentation to address class imbalance and to improve overall model performance. By integrating insights gained from preliminary investigations, a systematic approach to generating synthetic clinical narratives will be employed to bolster recognition of rare codes. The resulting models, once trained and evaluated, are expected to offer valuable contributions to automated ICD coding by delivering enhanced robustness and reliability.

\section{Gantt Chart of Planned Work Timeline}
\begin{figure}[H]
    \centering
    \begin{ganttchart}[
        hgrid,
        vgrid,
        x unit=0.8cm,
        y unit title=0.6cm,
        y unit chart=0.6cm,
        title height=1,
        bar/.style={fill=blue!50},
        bar height=0.5
    ]{1}{10}
        \gantttitle{Project Timeline: Sep 2024 -- Jun 2025}{10} \\
        \gantttitlelist{1,...,10}{1} \\
        \ganttbar{Phase 1: Literature Review}{1}{3} \\
        \ganttbar{Phase 2: Dataset Understanding}{2}{4} \\
        \ganttbar{Phase 3: Creating Co-occurring Rare Diseases}{4}{6} \\
        \ganttbar{Phase 4: Synthetic Dataset Generation}{6}{8} \\
        \ganttbar{Phase 5: Training and Evaluation}{8}{10} \\
    \end{ganttchart}
    \caption{Gantt Chart Reflecting Main Phases and Overlaps}
\end{figure}